\chapter{Methodology}
\label{chap:methodology}

\section{Victim Architectures}
\label{sec:architectures}

Three victim architectures are used across the attacks studied in this
project:

\begin{itemize}[leftmargin=1.6em]
  \item \textbf{LeNet-style CNN} - four convolutional layers followed by a
  fully-connected head. Two variants are used: the original DLG
  formulation (bias, sigmoid activations), used for DLG/iDLG and, under
  the same architecture, initial weights, and image, for Inverting
  Gradients, providing a common architecture for the three-attack
  comparison in Section~\ref{sec:dpsgd}; and a bias-free, LeakyReLU variant
  required by R-GAP (referred to here as LeNetRGAP), since R-GAP's
  per-layer equations (Section~\ref{sec:rgap}) assume an affine map with no
  additive bias and an invertible activation; its output layer is a single
  logit trained with binary logistic loss, since R-GAP's closed-form
  output-layer inversion is only defined for a scalar, sigmoid-activated
  loss.
  \item \textbf{ResNet-18} - used for Inverting Gradients, to test attack
  scalability on a deeper, residual architecture closer to what is
  deployed in practice.
  \item \textbf{FCN3} - a three-layer, bias-free, fully-connected network
  (1000, 100, 1 units, LeakyReLU throughout), used alongside LeNetRGAP for
  R-GAP as a point of comparison: it has no
  convolutional layers, and therefore no extra forward-pass constraints
  beyond the per-layer gradient equations themselves (Section~\ref{sec:rgap}).
\end{itemize}

\section{DLG and iDLG}
\label{sec:dlg}

DLG~\cite{zhu2019dlg} initializes a dummy sample $(\hat{x}, \hat{y})$ from
random noise and optimizes
\begin{equation}
\min_{\hat{x}, \hat{y}} \; \big\lVert g' - g \big\rVert_2^2,
\qquad g' = \nabla_w \, \ell(w_t; \hat{x}, \hat{y}),
\label{eq:dlg-objective}
\end{equation}
by gradient descent on $(\hat{x}, \hat{y})$ jointly. Because $g'$ is itself
a gradient, minimizing Equation~\eqref{eq:dlg-objective} requires
differentiating through $\nabla_w \ell$ a second time (``double
backpropagation''); both $\hat{x}$ and $\hat{y}$ are free, continuous
optimization variables, including the label, which DLG only relaxes to a
one-hot vector after optimization.

iDLG~\cite{zhao2020idlg} removes the label from the optimization. For
cross-entropy loss on a one-hot label $y$, the gradient of the loss with
respect to the pre-softmax logit $z_c$ of class $c$ is $\partial \ell/\partial z_c = p_c - \mathbb{1}[c = y]$,
where $p_c$ is the predicted probability. Since $p_c \in (0,1)$, this
quantity is negative only for the true class, so
\begin{equation}
\hat{y} = \arg\min_{c} \; \left(\frac{\partial \ell}{\partial z_c}\right),
\label{eq:idlg-label}
\end{equation}
recovers the ground-truth label in closed form from the sign pattern of the
last layer's gradient alone. Only $\hat{x}$ is then optimized against
Equation~\eqref{eq:dlg-objective}, which converges markedly faster and more
reliably than DLG's joint optimization.

\section{Inverting Gradients}
\label{sec:inverting-gradients}

Inverting Gradients~\cite{geiping2020inverting} replaces the $L_2$ distance
with a cosine-similarity cost and adds a total-variation image prior:
\begin{equation}
\min_{\hat{x}} \; \left(1 - \frac{\langle g', g \rangle}{\lVert g' \rVert \, \lVert g \rVert}\right)
+ \alpha \, \mathrm{TV}(\hat{x}), \qquad g' = \nabla_w \, \ell(w_t; \hat{x}, y),
\label{eq:ig-objective}
\end{equation}
where $\mathrm{TV}(\hat{x}) = \sum_{i,j} |\hat{x}_{i+1,j} - \hat{x}_{i,j}| + |\hat{x}_{i,j+1} - \hat{x}_{i,j}|$
is the total variation of the image - the sum of absolute differences
between horizontally and vertically neighboring pixels - which penalizes
non-smooth reconstructions. The label $y$ is recovered the same way as in
iDLG (Equation~\eqref{eq:idlg-label}), and $\hat{x}$ is constrained to the
valid image range $[0,1]^d$ by a box constraint. Two choices make
Equation~\eqref{eq:ig-objective} scale to deep networks where DLG's $L_2$
cost does not: the cosine-similarity term depends only on the
\emph{direction} of $g'$, not its magnitude, which is otherwise dominated
by whichever layer happens to have the largest gradient norm; and the
optimizer applies Adam updates using only the \emph{sign} of the computed
gradient, which prevents any single coordinate's magnitude from
controlling the step size. This project evaluates Equation~\eqref{eq:ig-objective}
on an untrained ResNet-18 victim model, and separately on the LeNet
architecture of Section~\ref{sec:architectures}, under the same weights
and image used for DLG/iDLG; the latter configuration is the one used in
the defense comparison of Section~\ref{sec:dpsgd}.

\section{R-GAP}
\label{sec:rgap}

R-GAP~\cite{zhu2021rgap} targets a feed-forward network of $L$ bias-free
layers with invertible elementwise activation $\sigma$: writing $a_0 = x$
for the input and, for $l = 1, \dots, L$,
\begin{equation}
z_l = W_l \, a_{l-1}, \qquad a_l = \sigma(z_l),
\label{eq:rgap-forward}
\end{equation}
the weight gradient at layer $l$ for a single sample ($B=1$) is the outer
product
\begin{equation}
\frac{\partial \ell}{\partial W_l} = \frac{\partial \ell}{\partial z_l} \, a_{l-1}^{\top}.
\label{eq:rgap-weight-grad}
\end{equation}
The leaked gradient $\partial \ell / \partial W_l$ is known for every
layer; $\partial \ell/\partial z_l$ is obtained recursively via the chain
rule, $\partial \ell/\partial z_l = \sigma'(z_l) \odot \big(W_{l+1}^{\top} \, \partial \ell/\partial z_{l+1}\big)$,
starting from the output layer, where $\partial \ell/\partial z_L$ is a
scalar equation solved numerically (the loss-to-logit derivative is a
monotonic function of the scalar logit for the binary logistic loss used
here). With $\partial \ell/\partial z_l$ known, Equation~\eqref{eq:rgap-weight-grad}
becomes a linear system for the unknown $a_{l-1}$, solved by least squares;
recursing this from $l = L$ down to $l = 1$ recovers $a_0 = x$ - the
private input itself - with \emph{no} optimization loop: each layer's
equation is solved exactly, subject to its own rank/invertibility
condition. Convolutional layers are handled by first re-expressing the
weight-sharing convolution as an equivalent dense (circulant) matrix
multiplication, so that Equation~\eqref{eq:rgap-weight-grad} applies
unchanged.

This project implements R-GAP on both victim architectures from
Section~\ref{sec:architectures} (LeNetRGAP and FCN3), and additionally runs
the entire forward/backward pass and every linear solve in float64 rather
than the default float32: because each layer's solution is an exact linear
system rather than one term in an averaged loss, small floating-point
rounding error does not cancel out as it would under gradient descent, and
is instead compounded across the $L$ recursive layers.

\section{DP-SGD as a Defense}
\label{sec:dpsgd}

DP-SGD~\cite{abadi2016dpsgd} is applied here, after the fact, as a
gradient-perturbation defense: before the leaked gradient $g$ reaches the
attacker, it is clipped to a global $L_2$ norm $C$ and perturbed with
Gaussian noise of multiplier $\sigma$,
\begin{equation}
\tilde{g} = \mathrm{clip}(g, C) + \mathcal{N}\!\big(0,\, (\sigma C)^2 I\big),
\qquad
\mathrm{clip}(g, C) = g \cdot \min\!\left(1, \frac{C}{\lVert g \rVert_2}\right).
\label{eq:dpsgd}
\end{equation}
The same perturbed gradient $\tilde{g}$ is substituted wherever $g$ appears
in Equations~\eqref{eq:dlg-objective}, \eqref{eq:ig-objective}, and the
R-GAP recursion of Section~\ref{sec:rgap}, allowing the three attacks to
be compared under an identical defense budget $(C, \sigma)$. Inverting
Gradients is evaluated under both victim configurations introduced in
Section~\ref{sec:inverting-gradients}: the ResNet-18 setup, and the LeNet
architecture shared with DLG/iDLG. This project sweeps $(C, \sigma)$ for
each attack and configuration; the resulting reconstruction quality is
reported in Chapter 5.
